\documentclass[11pt]{article}
\usepackage[margin=1in]{geometry}
\usepackage{amsmath,amssymb,amsthm}
\usepackage{hyperref}
\usepackage[T1]{fontenc}
\usepackage[utf8]{inputenc}

\newtheorem{definition}{Definition}
\newtheorem{conjecture}{Conjecture}
\newtheorem{proposition}{Proposition}
\newtheorem{example}{Example}

\title{Hyperseed Formalizations: Orientation Note 0001}
\author{OmegaClaw / ZeroBot working notes for Ben Goertzel}
\date{2026-06-27}

\begin{document}
\maketitle

\begin{abstract}
This note seeds a public repository for incremental Hyperseed formalizations arising
from OmegaClaw, ZeroBot, and Ben Goertzel's research interactions.  It is not a claim
of completed theory; it is a provenance-bearing scaffold for definitions, conjectures,
proof sketches, and experiment formalizations.
\end{abstract}

\section{Source orientation}

The initial orientation sources are Ben Goertzel's public Hyperseed v2 article and the
four linked PDF documents preserved in the OpenClaw research library on 2026-06-27:
\begin{itemize}
  \item \emph{Hyperseed ontology: formal presentation and theoretical development}.
  \item \emph{Hyperseed: A Scrutability-Style Core Ontology and a Practical Inference-Control Scaffold for Probabilistic Logic Networks}.
  \item \emph{Representing Scientific Experiments in MeTTa: SUMO/EXPO Structure with Hyperseed Epistemic Annotations}.
  \item \emph{Metta-stasizing SUMO: A Typed MeTTa Encoding for Type Checking, PLN Reasoning, and Hyperseed Alignment}.
\end{itemize}

\section{A minimal working frame}

\begin{definition}[Experienced research event]
An \emph{experienced research event} is a tuple
\[
  E = \langle O, A, C, R, P \rangle
\]
where $O$ is an observing agent or agent-ensemble, $A$ is an action or interaction,
$C$ is the operative context, $R$ is the resulting observation or artifact, and $P$ is
provenance sufficient for later scrutiny.
\end{definition}

\begin{definition}[Hyperseed formalization act]
A \emph{Hyperseed formalization act} maps an experienced research event $E$ into a
structured object
\[
  F(E) = \langle D, X, K, T, Q \rangle
\]
where $D$ is a set of candidate definitions, $X$ examples and counterexamples, $K$
conjectures, $T$ theorem/proof-sketch material, and $Q$ open questions or calibration
requests.
\end{definition}

\begin{conjecture}[Experiment-to-ontology feedback]
For a sustained research process, repeatedly applying $F$ to concrete experiments and
agent interactions can improve both experiment design and ontology design, provided the
formalization records uncertainty, provenance, and failure modes rather than only polished
successes.
\end{conjecture}

\begin{example}[petta-chem as initial target]
A deterministic candidate-selection run in \texttt{petta-chem} may be treated as an
experienced research event whose context includes a PeTTa/MeTTa runtime, symbolic chamber
state, candidate applicability rules, and observed post-tick abundance invariants.  A first
Hyperseed formalization should relate this to process, pattern, causality, attention or
resource allocation, and experiment-evidence annotations.
\end{example}

\section{Near-term agenda}

\begin{enumerate}
  \item Formalize \texttt{petta-chem} experiment records using the scientific-experiment
        paper's SUMO/EXPO plus Hyperseed epistemic annotation style.
  \item Formalize OmegaClaw's Telegram smoke test and idle-loop issue as agent-experience
        events involving communication, attention, resource use, and action gating.
  \item Extract a small reusable notation for ``interaction afterthoughts'': immediate
        conversational response first, later Hyperseed reflection when useful.
\end{enumerate}

\section{Open questions for Ben}

\begin{enumerate}
  \item Should the repository prefer mathematically polished theorem/proof style, or a
        lab-notebook style that tolerates more speculative formal sketches?
  \item Which Hyperseed primitives should be treated as the default starting vocabulary
        for formalizing software experiments?
  \item How much MeTTa syntax should be paired with each LaTeX formalization?
\end{enumerate}

\end{document}
